\documentclass{article}
\usepackage{lmodern}
\usepackage{amsmath}
\usepackage{listings}
\usepackage{fullpage}
\usepackage{parskip}
\usepackage{xcolor}
\lstset{
	basicstyle=\ttfamily,
	columns=fullflexible,
	frame=single,
	breaklines=true,
	postbreak=\mbox{\textcolor{red}{$\hookrightarrow$}\space},
}
\begin{document}
\title{Experiment `p\_13\_10\_recursive\_area\_polygon\_array' Results}
\author{\today}
\date{}
\maketitle
\textbf{Experiment outcome: }FAILURE\_VERIFY
\\\textbf{Bad responses: }1
\\\textbf{Responses containing}~\texttt{assume}~\textbf{: }0
\\\textbf{Responses containing}~\texttt{expect}~\textbf{: }0
\\\textbf{Resolution attempts: }4
\\\textbf{Hard fails (resolution): }1
\\\textbf{Soft fails (resolution): }0
\\\textbf{Verification attempts: }3
\section*{Problem Specification}
\textbf{Problem name: }p\_13\_10\_recursive\_area\_polygon\_array
\\\textbf{Natural language statement: }Write a recursive method that takes an array of points of a polygon and returns its area. Cut off a triangle and use the fact that a triangle with corners (x1, y1), (x2, y2), (x3, y3) has area (abs(x1y2 + x2y3 + x3y1 - y1x2 - y2x3 - y3x1)) / 2.
\\\textbf{Method signature: }p\_13\_10\_recursive\_area\_polygon\_array(points: array<(real, real)>) returns (area: real)
\subsection*{Ensures}
\begin{itemize}
\item \texttt{area == polygon\_area(points[..])}
\end{itemize}
\subsection*{Requires}
\begin{itemize}
\item \texttt{points.Length >= 3}
\end{itemize}
\subsection*{Functional Code Given}
\begin{lstlisting}
function polygon_area(points: seq<(real, real)>): real
  decreases |points|
  requires |points| >= 3
  ensures polygon_area(points) >= 0.0
{
  if |points| == 3 then 
    var (x1, y1) := points[0];
    var (x2, y2) := points[1];
    var (x3, y3) := points[2];
    var sign : real := x1 * (y2 - y3) + x2 * (y3 - y1) + x3 * (y1 - y2);
    if sign < 0.0 then -sign / 2.0 else sign / 2.0 as real 
  else 
    var (x1, y1) := points[0];
    var (x2, y2) := points[1];
    var (x3, y3) := points[2];
    var sign : real := x1 * (y2 - y3) + x2 * (y3 - y1) + x3 * (y1 - y2);
    var triangle_area := if sign < 0.0 then -sign / 2.0 else sign / 2.0 as real;
    triangle_area + polygon_area([points[0]] + points[2..])
}
\end{lstlisting}
\clearpage
\section*{GenAI interactions}
Below you will find all interactions between the `user' (program) and the `assistant' (OpenAI).
\subsection*{Program $\rightarrow$ GenAI}
\begin{lstlisting}
You are given the following task to perform in Dafny:

Write a recursive method that takes an array of points of a polygon and returns its area. Cut off a triangle and use the fact that a triangle with corners (x1, y1), (x2, y2), (x3, y3) has area (abs(x1y2 + x2y3 + x3y1 - y1x2 - y2x3 - y3x1)) / 2.

The signature should be:

p_13_10_recursive_area_polygon_array(points: array<(real, real)>) returns (area: real)

The method should respect the following contract:

requires points.Length >= 3, ensures area == polygon_area(points[..])

The contract uses the following dafny code:

function polygon_area(points: seq<(real, real)>): real
  decreases |points|
  requires |points| >= 3
  ensures polygon_area(points) >= 0.0
{
  if |points| == 3 then 
    var (x1, y1) := points[0];
    var (x2, y2) := points[1];
    var (x3, y3) := points[2];
    var sign : real := x1 * (y2 - y3) + x2 * (y3 - y1) + x3 * (y1 - y2);
    if sign < 0.0 then -sign / 2.0 else sign / 2.0 as real 
  else 
    var (x1, y1) := points[0];
    var (x2, y2) := points[1];
    var (x3, y3) := points[2];
    var sign : real := x1 * (y2 - y3) + x2 * (y3 - y1) + x3 * (y1 - y2);
    var triangle_area := if sign < 0.0 then -sign / 2.0 else sign / 2.0 as real;
    triangle_area + polygon_area([points[0]] + points[2..])
}

These function(s) must not be used in your implementation of the method.

Produce and show only the Dafny body of this method, including the curly braces that surround it. Do not show the signature nor contract. You must not use 'assume' nor 'expect' anywhere in your code.
\end{lstlisting}
\subsection*{GenAI $\rightarrow$ Program}
\textbf{System fingerprint: }null
\\\textbf{ID: }chatcmpl-DnAzkjLrO7L0dmnk718lviKV17ZGE
\\\textbf{Model: }gpt-5.5-2026-04-23
\\\textbf{Created at: }1780612708
\\\textbf{Finish reason: }stop
\\\textbf{Usage: }1058 tokens in, and 6415 tokens out
\begin{lstlisting}
{
  var (x1, y1) := points[0];
  var (x2, y2) := points[1];
  var (x3, y3) := points[2];
  var sign: real := x1 * (y2 - y3) + x2 * (y3 - y1) + x3 * (y1 - y2);
  var triangle_area: real := if sign < 0.0 then -sign / 2.0 else sign / 2.0;

  assert points[..][0] == points[0];
  assert points[..][1] == points[1];
  assert points[..][2] == points[2];

  if points.Length == 3 {
    area := triangle_area;
  } else {
    assert points.Length > 3;

    var rest := new (real, real)[points.Length - 1];
    rest[0] := points[0];

    var i := 2;
    while i < points.Length
      invariant 2 <= i <= points.Length
      invariant rest.Length == points.Length - 1
      invariant rest[0] == points[0]
      invariant forall j :: 1 <= j < i - 1 ==> rest[j] == points[j + 1]
    {
      rest[i - 1] := points[i];
      i := i + 1;
    }

    assert i == points.Length;
    assert forall j :: 1 <= j < rest.Length ==> rest[j] == points[j + 1];

    assert |rest[..]| == |[points[0]] + points[2..]|;
    assert forall j :: 0 <= j < |rest[..]| ==> rest[..][j] == ([points[0]] + points[2..])[j] by {
      assert rest[..][j] == rest[j];
      if j == 0 {
        assert rest[j] == points[0];
        assert ([points[0]] + points[2..])[j] == points[0];
      } else {
        assert 1 <= j < rest.Length;
        assert rest[j] == points[j + 1];
        assert 0 <= j - 1;
        assert j - 1 < |points[2..]|;
        assert ([points[0]] + points[2..])[j] == points[2..][j - 1];
        assert points[2..][j - 1] == points[2 + (j - 1)];
        assert 2 + (j - 1) == j + 1;
      }
    }
    assert rest[..] == [points[0]] + points[2..];
    assert points[..][2..] == points[2..];
    assert [points[..][0]] + points[..][2..] == [points[0]] + points[2..];

    var subArea := p_13_10_recursive_area_polygon_array(rest);
    area := triangle_area + subArea;
  }
}
\end{lstlisting}
\subsection*{Program $\rightarrow$ GenAI}
\begin{lstlisting}
When using dafny resolve, the below error is emitted and resolve fails:

problem_data/p_13_10_recursive_area_polygon_array/program.dfy(59,22): Error: unresolved identifier: j
   |
59 |       assert rest[..][j] == rest[j];
   |                       ^

problem_data/p_13_10_recursive_area_polygon_array/program.dfy(59,33): Error: unresolved identifier: j
   |
59 |       assert rest[..][j] == rest[j];
   |                                  ^

problem_data/p_13_10_recursive_area_polygon_array/program.dfy(60,9): Error: unresolved identifier: j
   |
60 |       if j == 0 {
   |          ^

problem_data/p_13_10_recursive_area_polygon_array/program.dfy(61,20): Error: unresolved identifier: j
   |
61 |         assert rest[j] == points[0];
   |                     ^

problem_data/p_13_10_recursive_area_polygon_array/program.dfy(62,43): Error: unresolved identifier: j
   |
62 |         assert ([points[0]] + points[2..])[j] == points[0];
   |                                            ^

problem_data/p_13_10_recursive_area_polygon_array/program.dfy(64,20): Error: unresolved identifier: j
   |
64 |         assert 1 <= j < rest.Length;
   |                     ^

problem_data/p_13_10_recursive_area_polygon_array/program.dfy(65,20): Error: unresolved identifier: j
   |
65 |         assert rest[j] == points[j + 1];
   |                     ^

problem_data/p_13_10_recursive_area_polygon_array/program.dfy(65,33): Error: unresolved identifier: j
   |
65 |         assert rest[j] == points[j + 1];
   |                                  ^

problem_data/p_13_10_recursive_area_polygon_array/program.dfy(66,20): Error: unresolved identifier: j
   |
66 |         assert 0 <= j - 1;
   |                     ^

problem_data/p_13_10_recursive_area_polygon_array/program.dfy(67,15): Error: unresolved identifier: j
   |
67 |         assert j - 1 < |points[2..]|;
   |                ^

problem_data/p_13_10_recursive_area_polygon_array/program.dfy(68,43): Error: unresolved identifier: j
   |
68 |         assert ([points[0]] + points[2..])[j] == points[2..][j - 1];
   |                                            ^

problem_data/p_13_10_recursive_area_polygon_array/program.dfy(68,61): Error: unresolved identifier: j
   |
68 |         assert ([points[0]] + points[2..])[j] == points[2..][j - 1];
   |                                                              ^

problem_data/p_13_10_recursive_area_polygon_array/program.dfy(69,27): Error: unresolved identifier: j
   |
69 |         assert points[2..][j - 1] == points[2 + (j - 1)];
   |                            ^

problem_data/p_13_10_recursive_area_polygon_array/program.dfy(69,49): Error: unresolved identifier: j
   |
69 |         assert points[2..][j - 1] == points[2 + (j - 1)];
   |                                                  ^

problem_data/p_13_10_recursive_area_polygon_array/program.dfy(70,20): Error: unresolved identifier: j
   |
70 |         assert 2 + (j - 1) == j + 1;
   |                     ^

problem_data/p_13_10_recursive_area_polygon_array/program.dfy(70,30): Error: unresolved identifier: j
   |
70 |         assert 2 + (j - 1) == j + 1;
   |                               ^

16 resolution/type errors detected in program.dfy

Correct the error in code by altering only the method body. Produce and show only the Dafny body, including the curly braces that surround it. Do not show the signature nor contract. You must not use 'assume' nor 'expect' anywhere in your code.
\end{lstlisting}
\subsection*{GenAI $\rightarrow$ Program}
\textbf{System fingerprint: }null
\\\textbf{ID: }chatcmpl-DnB1pKczgq6I4PQ5Yr3vTshemA5lL
\\\textbf{Model: }gpt-5.5-2026-04-23
\\\textbf{Created at: }1780612837
\\\textbf{Finish reason: }stop
\\\textbf{Usage: }3047 tokens in, and 1756 tokens out
\begin{lstlisting}
{
  var (x1, y1) := points[0];
  var (x2, y2) := points[1];
  var (x3, y3) := points[2];
  var sign: real := x1 * (y2 - y3) + x2 * (y3 - y1) + x3 * (y1 - y2);
  var triangle_area: real := if sign < 0.0 then -sign / 2.0 else sign / 2.0;

  assert points[..][0] == points[0];
  assert points[..][1] == points[1];
  assert points[..][2] == points[2];

  if points.Length == 3 {
    area := triangle_area;
  } else {
    assert points.Length > 3;

    var rest := new (real, real)[points.Length - 1];
    rest[0] := points[0];

    var i := 2;
    while i < points.Length
      invariant 2 <= i <= points.Length
      invariant rest.Length == points.Length - 1
      invariant rest[0] == points[0]
      invariant forall j :: 1 <= j < i - 1 ==> rest[j] == points[j + 1]
    {
      rest[i - 1] := points[i];
      i := i + 1;
    }

    assert i == points.Length;
    assert forall j :: 1 <= j < rest.Length ==> rest[j] == points[j + 1];

    assert |rest[..]| == |[points[0]] + points[2..]|;
    forall j | 0 <= j < |rest[..]|
      ensures rest[..][j] == ([points[0]] + points[2..])[j]
    {
      assert rest[..][j] == rest[j];
      if j == 0 {
        assert rest[j] == points[0];
        assert ([points[0]] + points[2..])[j] == points[0];
      } else {
        assert 1 <= j < rest.Length;
        assert rest[j] == points[j + 1];
        assert 0 <= j - 1;
        assert j - 1 < |points[2..]|;
        assert ([points[0]] + points[2..])[j] == points[2..][j - 1];
        assert points[2..][j - 1] == points[2 + (j - 1)];
        assert 2 + (j - 1) == j + 1;
      }
    }
    assert rest[..] == [points[0]] + points[2..];
    assert points[..][2..] == points[2..];
    assert [points[..][0]] + points[..][2..] == [points[0]] + points[2..];

    var subArea := p_13_10_recursive_area_polygon_array(rest);
    area := triangle_area + subArea;
  }
}
\end{lstlisting}
\subsection*{Program $\rightarrow$ GenAI}
\begin{lstlisting}
When using dafny verify, the below error is emitted and verify fails:

problem_data/p_13_10_recursive_area_polygon_array/program.dfy(79,55): Error: cannot prove termination; try supplying a decreases clause
   |
79 |     var subArea := p_13_10_recursive_area_polygon_array(rest);
   |                                                        ^


Dafny program verifier finished with 2 verified, 1 error

Correct the error in code by altering only the method body. Produce and show only the Dafny body, including the curly braces that surround it. Do not show the signature nor contract. You must not use 'assume' nor 'expect' anywhere in your code.
\end{lstlisting}
\subsection*{GenAI $\rightarrow$ Program}
\textbf{System fingerprint: }null
\\\textbf{ID: }chatcmpl-DnB2ZECh8tpvLJjrx6MBGjWz91s5Y
\\\textbf{Model: }gpt-5.5-2026-04-23
\\\textbf{Created at: }1780612883
\\\textbf{Finish reason: }stop
\\\textbf{Usage: }3166 tokens in, and 5696 tokens out
\begin{lstlisting}
{
  var total: real := 0.0;
  var k: int := 1;

  assert points.Length >= 3;
  assert 1 <= points.Length - 2;

  assert |[points[0]] + points[1..]| == |points[..]|;
  forall j | 0 <= j < |points[..]|
    ensures ([points[0]] + points[1..])[j] == points[..][j]
  {
    if j == 0 {
      assert ([points[0]] + points[1..])[j] == points[0];
      assert points[..][j] == points[0];
    } else {
      assert 0 <= j - 1 < |points[1..]|;
      assert ([points[0]] + points[1..])[j] == points[1..][j - 1];
      assert points[1..][j - 1] == points[j];
      assert points[..][j] == points[j];
    }
  }
  assert [points[0]] + points[1..] == points[..];

  while k < points.Length - 2
    invariant 1 <= k <= points.Length - 2
    invariant total + polygon_area([points[0]] + points[k..]) == polygon_area(points[..])
    decreases points.Length - 2 - k
  {
    var oldTail := [points[0]] + points[k..];
    assert |oldTail| > 3;
    assert oldTail[0] == points[0];
    assert oldTail[1] == points[k];
    assert oldTail[2] == points[k + 1];

    var (x1, y1) := points[0];
    var (x2, y2) := points[k];
    var (x3, y3) := points[k + 1];
    var sign: real := x1 * (y2 - y3) + x2 * (y3 - y1) + x3 * (y1 - y2);
    var tri: real := if sign < 0.0 then -sign / 2.0 else sign / 2.0;

    var newTail := [points[0]] + points[k + 1..];

    assert |oldTail[2..]| == |points[k + 1..]|;
    forall j | 0 <= j < |oldTail[2..]|
      ensures oldTail[2..][j] == points[k + 1..][j]
    {
      assert oldTail[2..][j] == oldTail[j + 2];
      assert oldTail[j + 2] == points[k..][j + 1];
      assert points[k..][j + 1] == points[k + j + 1];
      assert points[k + 1..][j] == points[k + 1 + j];
    }
    assert oldTail[2..] == points[k + 1..];
    assert [oldTail[0]] + oldTail[2..] == newTail;
    assert polygon_area(oldTail) == tri + polygon_area(newTail);
    assert total + polygon_area(oldTail) == polygon_area(points[..]);
    assert total + tri + polygon_area(newTail) == polygon_area(points[..]);

    total := total + tri;
    k := k + 1;
    assert [points[0]] + points[k..] == newTail;
    assert total + polygon_area([points[0]] + points[k..]) == polygon_area(points[..]);
  }

  assert k == points.Length - 2;
  var tail := [points[0]] + points[k..];
  assert |tail| == 3;
  assert tail[0] == points[0];
  assert tail[1] == points[k];
  assert tail[2] == points[k + 1];

  var (x1, y1) := points[0];
  var (x2, y2) := points[k];
  var (x3, y3) := points[k + 1];
  var sign: real := x1 * (y2 - y3) + x2 * (y3 - y1) + x3 * (y1 - y2);
  var tri: real := if sign < 0.0 then -sign / 2.0 else sign / 2.0;

  assert polygon_area(tail) == tri;
  area := total + tri;
  assert area == polygon_area(points[..]);
}
\end{lstlisting}
\subsection*{Program $\rightarrow$ GenAI}
\begin{lstlisting}
When using dafny verify, the below error is emitted and verify fails:

problem_data/p_13_10_recursive_area_polygon_array/program.dfy(77,33): Error: assertion might not hold
   |
77 |     assert polygon_area(oldTail) == tri + polygon_area(newTail);
   |                                  ^^


Dafny program verifier finished with 2 verified, 1 error

Correct the error in code by altering only the method body. Produce and show only the Dafny body, including the curly braces that surround it. Do not show the signature nor contract. You must not use 'assume' nor 'expect' anywhere in your code.
\end{lstlisting}
\subsection*{GenAI $\rightarrow$ Program}
\textbf{System fingerprint: }null
\\\textbf{ID: }chatcmpl-DnB4TTzmhdk2basZsae3IQbIephXz
\\\textbf{Model: }gpt-5.5-2026-04-23
\\\textbf{Created at: }1780613001
\\\textbf{Finish reason: }stop
\\\textbf{Usage: }4375 tokens in, and 3205 tokens out
\begin{lstlisting}
{
  var total: real := 0.0;
  var k: int := 1;

  assert |[points[0]] + points[1..]| == |points[..]|;
  forall j | 0 <= j < |points[..]|
    ensures ([points[0]] + points[1..])[j] == points[..][j]
  {
    if j == 0 {
      assert ([points[0]] + points[1..])[j] == points[0];
      assert points[..][j] == points[0];
    } else {
      assert 0 <= j - 1 < |points[1..]|;
      assert ([points[0]] + points[1..])[j] == points[1..][j - 1];
      assert points[1..][j - 1] == points[j];
      assert points[..][j] == points[j];
    }
  }
  assert [points[0]] + points[1..] == points[..];

  while k < points.Length - 2
    invariant 1 <= k <= points.Length - 2
    invariant total + polygon_area([points[0]] + points[k..]) == polygon_area(points[..])
    decreases points.Length - 2 - k
  {
    var oldTail := [points[0]] + points[k..];
    assert |oldTail| > 3;
    assert oldTail[0] == points[0];
    assert oldTail[1] == points[k];
    assert oldTail[2] == points[k + 1];

    var (sx1, sy1) := oldTail[0];
    var (sx2, sy2) := oldTail[1];
    var (sx3, sy3) := oldTail[2];
    var ssign: real := sx1 * (sy2 - sy3) + sx2 * (sy3 - sy1) + sx3 * (sy1 - sy2);
    var stri: real := if ssign < 0.0 then -ssign / 2.0 else ssign / 2.0;

    var newTail := [points[0]] + points[k + 1..];

    assert |oldTail[2..]| == |points[k + 1..]|;
    forall j | 0 <= j < |oldTail[2..]|
      ensures oldTail[2..][j] == points[k + 1..][j]
    {
      assert oldTail[2..][j] == oldTail[j + 2];
      assert 0 < j + 2 < |oldTail|;
      assert oldTail[j + 2] == points[k..][j + 1];
      assert 0 <= j + 1 < |points[k..]|;
      assert points[k..][j + 1] == points[k + j + 1];
      assert points[k + 1..][j] == points[k + 1 + j];
      assert k + j + 1 == k + 1 + j;
    }
    assert oldTail[2..] == points[k + 1..];
    assert [oldTail[0]] + oldTail[2..] == newTail;

    assert polygon_area(oldTail) == stri + polygon_area([oldTail[0]] + oldTail[2..]);
    assert polygon_area(oldTail) == stri + polygon_area(newTail);
    assert total + polygon_area(oldTail) == polygon_area(points[..]);
    assert total + stri + polygon_area(newTail) == polygon_area(points[..]);

    total := total + stri;
    k := k + 1;
    assert [points[0]] + points[k..] == newTail;
    assert total + polygon_area([points[0]] + points[k..]) == polygon_area(points[..]);
  }

  assert k == points.Length - 2;
  var tail := [points[0]] + points[k..];
  assert |tail| == 3;
  assert tail[0] == points[0];
  assert tail[1] == points[k];
  assert tail[2] == points[k + 1];

  var (tx1, ty1) := tail[0];
  var (tx2, ty2) := tail[1];
  var (tx3, ty3) := tail[2];
  var tsign: real := tx1 * (ty2 - ty3) + tx2 * (ty3 - ty1) + tx3 * (ty1 - ty2);
  var ttri: real := if tsign < 0.0 then -tsign / 2.0 else tsign / 2.0;

  assert polygon_area(tail) == ttri;
  area := total + ttri;
  assert area == polygon_area(points[..]);
}
\end{lstlisting}
\subsection*{Program $\rightarrow$ GenAI}
\begin{lstlisting}
When using dafny verify, the below error is emitted and verify fails:

problem_data/p_13_10_recursive_area_polygon_array/program.dfy(78,33): Error: assertion might not hold
   |
78 |     assert polygon_area(oldTail) == stri + polygon_area([oldTail[0]] + oldTail[2..]);
   |                                  ^^

problem_data/p_13_10_recursive_area_polygon_array/program.dfy(102,28): Error: assertion might not hold
    |
102 |   assert polygon_area(tail) == ttri;
    |                             ^^


Dafny program verifier finished with 2 verified, 2 errors

Correct the error in code by altering only the method body. Produce and show only the Dafny body, including the curly braces that surround it. Do not show the signature nor contract. You must not use 'assume' nor 'expect' anywhere in your code.
\end{lstlisting}
\subsection*{GenAI $\rightarrow$ Program}
\textbf{System fingerprint: }null
\\\textbf{ID: }chatcmpl-DnB5ddhxH9BmRMNruiopfcjQLbkXS
\\\textbf{Model: }gpt-5.5-2026-04-23
\\\textbf{Created at: }1780613073
\\\textbf{Finish reason: }length
\\\textbf{Usage: }5775 tokens in, and 1282 tokens out
\\\textcolor{red}{\textbf{The below is a raw response from GenAI that could not be parsed. It is likely that GenAI got 'stuck' in a loop and hit the pre-set token limit (i.e., finish reason is 'length').}}
\begin{lstlisting}

\end{lstlisting}
\section*{Final Program}
\begin{lstlisting}
function polygon_area(points: seq<(real, real)>): real
  decreases |points|
  requires |points| >= 3
  ensures polygon_area(points) >= 0.0
{
  if |points| == 3 then 
    var (x1, y1) := points[0];
    var (x2, y2) := points[1];
    var (x3, y3) := points[2];
    var sign : real := x1 * (y2 - y3) + x2 * (y3 - y1) + x3 * (y1 - y2);
    if sign < 0.0 then -sign / 2.0 else sign / 2.0 as real 
  else 
    var (x1, y1) := points[0];
    var (x2, y2) := points[1];
    var (x3, y3) := points[2];
    var sign : real := x1 * (y2 - y3) + x2 * (y3 - y1) + x3 * (y1 - y2);
    var triangle_area := if sign < 0.0 then -sign / 2.0 else sign / 2.0 as real;
    triangle_area + polygon_area([points[0]] + points[2..])
}

method p_13_10_recursive_area_polygon_array(points: array<(real, real)>) returns (area: real)
	requires points.Length >= 3
	ensures area == polygon_area(points[..])
{
  var total: real := 0.0;
  var k: int := 1;

  assert |[points[0]] + points[1..]| == |points[..]|;
  forall j | 0 <= j < |points[..]|
    ensures ([points[0]] + points[1..])[j] == points[..][j]
  {
    if j == 0 {
      assert ([points[0]] + points[1..])[j] == points[0];
      assert points[..][j] == points[0];
    } else {
      assert 0 <= j - 1 < |points[1..]|;
      assert ([points[0]] + points[1..])[j] == points[1..][j - 1];
      assert points[1..][j - 1] == points[j];
      assert points[..][j] == points[j];
    }
  }
  assert [points[0]] + points[1..] == points[..];

  while k < points.Length - 2
    invariant 1 <= k <= points.Length - 2
    invariant total + polygon_area([points[0]] + points[k..]) == polygon_area(points[..])
    decreases points.Length - 2 - k
  {
    var oldTail := [points[0]] + points[k..];
    assert |oldTail| > 3;
    assert oldTail[0] == points[0];
    assert oldTail[1] == points[k];
    assert oldTail[2] == points[k + 1];

    var (sx1, sy1) := oldTail[0];
    var (sx2, sy2) := oldTail[1];
    var (sx3, sy3) := oldTail[2];
    var ssign: real := sx1 * (sy2 - sy3) + sx2 * (sy3 - sy1) + sx3 * (sy1 - sy2);
    var stri: real := if ssign < 0.0 then -ssign / 2.0 else ssign / 2.0;

    var newTail := [points[0]] + points[k + 1..];

    assert |oldTail[2..]| == |points[k + 1..]|;
    forall j | 0 <= j < |oldTail[2..]|
      ensures oldTail[2..][j] == points[k + 1..][j]
    {
      assert oldTail[2..][j] == oldTail[j + 2];
      assert 0 < j + 2 < |oldTail|;
      assert oldTail[j + 2] == points[k..][j + 1];
      assert 0 <= j + 1 < |points[k..]|;
      assert points[k..][j + 1] == points[k + j + 1];
      assert points[k + 1..][j] == points[k + 1 + j];
      assert k + j + 1 == k + 1 + j;
    }
    assert oldTail[2..] == points[k + 1..];
    assert [oldTail[0]] + oldTail[2..] == newTail;

    assert polygon_area(oldTail) == stri + polygon_area([oldTail[0]] + oldTail[2..]);
    assert polygon_area(oldTail) == stri + polygon_area(newTail);
    assert total + polygon_area(oldTail) == polygon_area(points[..]);
    assert total + stri + polygon_area(newTail) == polygon_area(points[..]);

    total := total + stri;
    k := k + 1;
    assert [points[0]] + points[k..] == newTail;
    assert total + polygon_area([points[0]] + points[k..]) == polygon_area(points[..]);
  }

  assert k == points.Length - 2;
  var tail := [points[0]] + points[k..];
  assert |tail| == 3;
  assert tail[0] == points[0];
  assert tail[1] == points[k];
  assert tail[2] == points[k + 1];

  var (tx1, ty1) := tail[0];
  var (tx2, ty2) := tail[1];
  var (tx3, ty3) := tail[2];
  var tsign: real := tx1 * (ty2 - ty3) + tx2 * (ty3 - ty1) + tx3 * (ty1 - ty2);
  var ttri: real := if tsign < 0.0 then -tsign / 2.0 else tsign / 2.0;

  assert polygon_area(tail) == ttri;
  area := total + ttri;
  assert area == polygon_area(points[..]);
}
\end{lstlisting}
\section*{Total Token Usage}
\textbf{Input tokens: }17421
\\\textbf{Output tokens: }18354
\\\textbf{Reasoning tokens: }14537
\\\textbf{Sum of `total tokens': }35775
\section*{Experiment Timings}
\textbf{Overall Experiment} started at 1780612707574, ended at 1780613102756, lasting 395182ms (395.18 seconds)
\\\textbf{Iteration \#4} started at 1780613000775, ended at 1780613073012, lasting 72237ms (72.24 seconds)
\\\textbf{Iteration \#5} started at 1780613073013, ended at 1780613102592, lasting 29579ms (29.58 seconds)
\\\textbf{Iteration \#1} started at 1780612707575, ended at 1780612836781, lasting 129206ms (129.21 seconds)
\\\textbf{Iteration \#2} started at 1780612836782, ended at 1780612883522, lasting 46740ms (46.74 seconds)
\\\textbf{Iteration \#3} started at 1780612883522, ended at 1780613000775, lasting 117253ms (117.25 seconds)
\end{document}
